\cotefile{34}
\scanpage{180}
\mspage{15}
% archive sheet number: 179
% block: top, line 1 | text
3.2. Démonstration de \unsure{1.4|1.9} à l'aide de 3.1.

% block: upper third, lines 2-7 | text
Comme $U_0$ est \unsure{l'ens|l'ouv|l'un} des $t_0\in B_0$ \unsure{dont} la fibre $\overline X_{t_0}$ est lisse, on sait que $f$ applique $U$ dans $U_0$; donc, \deleted{\unsure{l'inertie est}} \deleted{$\unsure{\pi_1}$} en vertu du th. du changement de base pour les $\pi_1^{(p)'}$ (SGA 1 $\overline{\underline{\mathrm{XIII}}}$\ldots) on est ramené à prouver le th., \deleted{\unsure{dans}} dans la situation \unsure{avec indices 0}. Or, en vertu d'une \inserted{variante du} lemme de Abhyankar (SGA 1 $\overline{\underline{\mathrm{XIII}}}$\ldots), le groupe fondamental de $U_0$ est \deleted{abélien premier} abélien d'ordre premier à $p$.

% block: middle third, lines 8-11 | text
Cela achève donc la démonstration de 1.4 dans le cas \deleted{De façon précise, il est isomorphe} de la dim relative 1, \inserted{de 1.4 et de 1.3} donc dans le cas général, grâce aux réductions du n° 2.

% block: middle to lower third, lines 12-18 | text
\uline{Remarque} 3.2.1. De façon précise, la variante indiquée du lemme de Abhyankar affirme que l'on trouve un revêtement universel de $U$ en prenant \inserted{la limite} des revêtements de $U$ induits par les revêtements kummeriens $S[T_1,\ldots,T_{\unsure{\rho|r|p}}]/(T_1^m,\ldots,T_\rho^m)$, \uline{avec $m$} \uline{premier à} $p$. Donc on trouve un isom canonique

% block: lower third, line 19 | display
\[
\pi_1(U_0,\unsure{\zeta|\xi})\simeq \deleted{\unsure{\mathbb Z}}\mathbb Z'_{\deleted{\unsure{\ell}}}(1)^I,
\]

% block: lower third, lines 20-22 | text
$I$ étant l'ens d'indices pour \deleted{\illegible{1}} les pts \deleted{\unsure{singuliers}}\inserted{de non} lissité de $X_0$, et $\mathbb Z'_{\deleted{\unsure{\ell}}}(1)$ est le produit des $\mathbb Z_\ell(1)$ ($\ell$ nb premier, \inserted{avec} $\ell\ne p$), $\leftarrow\mathbb Z_\ell(1)=\varprojlim_m\mu_{\ell^m}(k)$.

% block: bottom, lines 23-25 | text
On peut donc dire que l'action de $\pi_1(U_0)$ est déterminée essentiellement par \deleted{un}\inserted{la donnée} \deleted{\unsure{certain}} d'\unsure{automorphismes} \inserted{$h_i$} de $\pi_1^{p'}(X_{\overline{\unsure{K|\eta|s}}})$, \deleted{\unsure{définis à}}\inserted{\unsure{couplés|accouplés}} (i.e. le choix d'un

\scanpage{181}
\mspage{16}
% archive sheet number: 180
% block: top, line 1 | text
générateur topologique de $\mathbb Z'(1)$) \deleted{\unsure{pour les}}\inserted{aux} \unsure{différents} $x_i$.

% block: upper third, lines 2-6 | text
\deleted{\unsure{La structure de}} Comme on \unsure{connait|connaît} la structure de $\pi_1^{\unsure{(p)'|p'}}(X_{\overline K},\unsure{\zeta|\xi})$ \inserted{en terme du genre de $X_{\overline K}$ et du nbre de pts \unsure{enlevés}} (SGA 1 $\overline{\underline{\mathrm{XIII}}}$) \deleted{\unsure{cela se lit}} on voit que le pb de déterminer la forme des $h_i$ devient assez concret. Cela \deleted{\unsure{pourra être à cet}} sera fait par Deligne \uline{par voie transcendante} dans un exposé ultérieur.

% block: upper to middle third, lines 7-10 | text
Lorsque $S$ est lui-\unsure{$=$} un trait, donc \deleted{$\unsure{\pi_1^{\prime}}$} $I_t=\pi_1(\eta,\overline\eta)_t\simeq\mathbb Z'_{\deleted{\unsure{\ell}}}(1)$, alors l'action extérieure de $I_t$ sur $\pi_1^{(p)'}(X_{\overline K},\zeta)$ est connue quand on \unsure{connait|connaît} l'hom induit $\deleted{\mathbb Z'(1)\mathbin{-}}\pi_1(\eta,\overline\eta)_{\mathrm{tam}}\longrightarrow\pi_1(U_0,\overline\eta)$ i.e. un \unsure{l'hom}

% block: middle third, lines 11-12 | display
\[
\begin{aligned}
&\deleted{\mathbb Z'(1)\mathbin{-}}\\
u:\ &\mathbb Z'(1)\longrightarrow\mathbb Z'(1)^I
\end{aligned}
\]

% block: middle third, lines 13-15 | text
\deleted{dont la connaissance \unsure{équivaut} d'ailleurs, comme on \unsure{voit aussitôt, à la}} qui est d'ailleurs donné par

% block: middle third, line 16 | display
\[
u(\lambda)=\deleted{\unsure{\sum}}(n_i\lambda)\deleted{\unsure{\text{avec}}}
\]

% block: middle to lower third, line 17 | display
\[
\text{avec}\qquad n_i=\operatorname{long}(V/\unsure{\mathfrak a_i|\mathfrak m_i}V).
\]

% block: lower third, lines 18-22 | text
où $\mathfrak a_i$ est l'idéal définissant le diviseur $D_i$\inserted{de $S$}. (L'entier $n_i$ mesure, \deleted{\unsure{en un sens}} heuristiquement, la vitesse avec laquelle \unsure{disparait|disparaît} la singularité $x_i$ en passant de la fibre spéciale à la fibre \unsure{générale|générique}.) Donc le générateur choisi de $\mathbb Z'(1)$ agit \inserted{extérieurement} sur $\pi_1^{p'}(X_{\overline K};\zeta)$ par

% block: bottom, line 23 | display
\[
h=h_1^{n_1}\cdots h_{\unsure{\rho|r|p}}^{n_\rho}
\]

\scanpage{183}
\mspage{II.31}
% archive sheet number: 182
% block: upper left, paired title strokes | drawing
\margin{\figure{\begin{tikzpicture}[x=1mm,y=1mm]
\draw (0,0)--(0,17);\draw (4,-2)--(4,11);
\end{tikzpicture}}}
% block: top, lines 1-3 | text
\uline{Appendice.}

6. \uline{Groupes et Algèbres de Lie \inserted{analytiques \unsure{ou} algébriques \deleted{de}} associés aux représentations $\ell$-adiques de groupes profinis}

% block: upper third, line 4 | text
On se fixe un nb premier $\ell$.

% block: upper third, lines 5-10 | text
6.1. Rappelons que sur un groupe topologique \unsure{admettant} une structure de groupe analytique sur $\mathbb Q_\ell$, cette dernière est unique (et si le groupe n'est pas discret, $\ell$ est uniquement déterminé); \deleted{\unsure{et tout}} \unsure{tout} \inserted{\unsure{tout}} hom. \uline{continu} de groupes analytiques sur $\mathbb Q_\ell$ est automatiquement continu. Enfin, un s-groupe \uline{fermé} d'un groupe analytique (sur $\mathbb Q_\ell$) est encore analytique.

% block: middle third, lines 11-16 | text
6.2. Soit \deleted{$\unsure{G}$} $\pi$ un groupe profini. Si $G=\pi/V$ et $G'=\pi/V'$ sont analytiques sur $\mathbb Q_\ell$, alors $G''=G/V\cap V'$ aussi, car il \deleted{\unsure{s'injecte}} est isomorphe à un sous-groupe fermé de $G\times G'$. Donc les \deleted{$G/V$} quotients \uline{analytiques} de $G$ forment un système projectif de groupes analytiques

% block: middle third, line 17 | display
\[
\deleted{\unsure{\mathscr G_\ell\pi}}\deleted{\unsure{\mathscr L_\ell(\pi)|\mathscr G_\ell(\pi)}}\inserted{\pi^{\mathrm{an}}}=(G_\alpha)
\]

% block: middle to lower third, lines 18-21 | text
qu'on appellera ``syst. projectif de Lie'' associé à $\pi$. \unsure{Noter|Notons} que si on considère $\pi$ comme \deleted{\unsure{pro-groupe}} \deleted{\unsure{profini}} syst. projectif de groupes \deleted{\unsure{analytiques}}\inserted{finis, considérés} comme groupes analytiques, on a un hom évident

% block: lower third, line 22 | display
\[
\deleted{\unsure{\mathscr L_\ell}}\pi^{\mathrm{an}}\mathbin{-}\pi.
\]

% block: lower third, lines 23-27 | text
6.3. Les représentations $\ell$-adiques de $\pi$ (sur $\mathbb Z_\ell$, ou sur $\mathbb Q_\ell$) ci-dessus s'interprètent alors comme représentations \deleted{\unsure{topologiques}}\inserted{\unsure{continus|continues} (analytiques)} de $\mathscr L_\ell(\pi)$, i.e. la catégorie des modules $\ell$-adiques sur $\pi$ est la \deleted{limite}\inserted{lim} des catégories analogues

\scanpage{184}
\mspage{II.32}
% archive sheet number: 183
% block: top, lines 1-2 | text
formées avec les $G_\alpha$, qu'on peut considérer comme modules à action analytique.

% block: upper third, lines 3-4 | text
6.4. Le syst. projectif des $G_\alpha$ définit un système projectif d'Algèbres de Lie $\mathfrak g_\alpha$ sur $\mathbb Q_\ell$

% block: upper third, line 5 | display
\[
\deleted{\unsure{\mathfrak g(\pi)=\mathfrak g_\ell(\pi)=(\mathfrak g_\alpha)}}\qquad
\mathfrak g=\operatorname{Lie}^{\mathrm{an}}(\pi)=(\mathfrak g_\alpha)
\]

% block: upper third, line 6 | text
et on a un foncteur

% block: upper to middle third, lines 7-11 | diagram
\diagram{\begin{tikzcd}[ampersand replacement=\&,column sep=large]
\& \deleted{\operatorname{Mod}(\unsure{G_\alpha})} \& \\
\operatorname{Mod}_{\mathbb Q_\ell}(\pi) \arrow[ur,no head,"{\unsure{\simeq}}"] \arrow[rr] \arrow[d]
\& \& \operatorname{Mod}\mathfrak g_{\deleted{\unsure{\ell}}} \arrow[d,no head] \\
\begin{gathered}\text{catégorie des}\\\text{\unsure{les modules|modules} $\ell$-adiques}\\\text{sur $\pi$ (sur $\mathbb Q_\ell$)}\end{gathered}
\& \& \begin{gathered}\varinjlim\text{ des}\\\text{catégories de}\\\text{modules }\deleted{\text{sur les}}\mathfrak g_\alpha\\\text{de dim finie}\\\text{sur }\mathbb Q_\ell\end{gathered}
\end{tikzcd}}

% block: middle third, lines 12-16 | text
Ce foncteur est fidèle, mais en général \uline{pas} pleinement fidèle: On a perdu en passant de \inserted{$\unsure{\pi^{\mathrm{an}}}$} $G$ à $\mathfrak g$ (\unsure{intuitivement}, on perd ``des composantes connexes). Mais soit $(U_i)$ le syst. fond. des \unsure{s-groupes|sous groupes} ouverts de $\pi$. Alors l'inclusion $U_i\mathbin{-}\pi$ définit un \uline{isom} de pro-Lie

% block: lower third, line 17 | display
\[
\mathfrak g(U_i)\simeq\mathfrak g(\pi)
\]

% block: lower third, line 18 | text
d'où un foncteur

% block: lower third, line 19 | display
\[
\operatorname{Mod}_{\mathbb Q_\ell}(U_i)\longrightarrow\operatorname{Mod}\mathfrak g
\]

% block: lower third, line 20 | text
et par passage à la limite

% block: lower third, line 21 | display
\[
\boxed{\varinjlim_i\operatorname{Mod}_{\mathbb Q_\ell}(U_i)\longrightarrow\operatorname{Mod}\mathfrak g}
\]

% block: bottom, line 22 | text
Cette flèche est une équivalence de \uline{catégories}

\scanpage{185}
\mspage{II.33}
% archive sheet number: 184
% block: top, lines 1-2 | text
Cela se ramène formellement \deleted{au} \deleted{\unsure{cas où $\pi$ est un groupe analytiquement} \illegible{2}}\inserted{\unsure{au thm bien connu}}

% block: upper third, lines 3-5 | text
\uline{\unsure{Lemme|Théorème}} Soit $G$ un groupe analytique, $\mathfrak g$ son alg. de Lie, qui est aussi l'alg. de Lie de ses sous-groupes ouverts \inserted{$U_i$}. Alors le foncteur

% block: upper third, line 6 | display
\[
\unsure{\varinjlim_i|\underset{i}{\underline{\lim}}}\operatorname{Mod}_{\mathbb Q_\ell}U_i\longrightarrow\operatorname{Mod}_{\deleted{\unsure{\mathbb Q_\ell}}}\mathfrak g
\]

% block: upper third, line 7 | text
est une équivalence de catégories.

% block: middle third, lines 8-13 | text
6.5. On fera attention \deleted{\unsure{qu'}}\inserted{que même si $\pi=\widehat{\mathbb Z}$,} on ne récupère pas $\pi$ (ou les \unsure{proobjets} $G$ \inserted{\unsure{qu'on vient associer}} (ni même le pro-Lie $\mathfrak g$) par la connaissance de la catégorie $\operatorname{Mod}_{\mathbb Q_\ell}(\pi)$, même munie de sa structure de produit tensoriel! on ne récupère que des enveloppes \uline{algébriques}. [\unsure{On montrera} cependant si \deleted{$\unsure{G}$} les $\mathfrak g_\alpha$ sont semi simples\ldots].

% block: lower left margin beside construction | margin note
% slanted
\margin{\fbox{\parbox{25mm}{Attention, il est essentiel ici de travailler sur $\mathbb Q_\ell$, non sur $\mathbb Z_\ell$; \unsure{le faire réserve des surprises}!}}}

% block: middle to lower third, lines 14-18 | text
\deleted{6.6} Je vais construire un groupe pro-\uline{algébrique} \inserted{\uline{affine}} canonique \inserted{$\pi^{\mathrm{alg}}$} sur \textcircled{$\mathbb Q_\ell$}, dont les représentations algébriques sur \deleted{$\unsure{\pi}$} $\mathbb Q_\ell$ vont définir la catégorie $\operatorname{Mod}_{\mathbb Q_\ell}(\pi)$ à équivalence près, en respectant la structure tensorielle. Pour ceci, on considère, pour tt $E\in\operatorname{Ob}\operatorname{Mod}_{\mathbb Q_\ell}(\pi)$, le groupe algébrique engendré dans $\underline{\operatorname{Aut}}_{\mathbb Q_\ell}(E)$ par l'image de $\pi$, soit $\widetilde G_E$. S'il existe un hom \uline{surjectif} $E\mathbin{-}F$, il y a un hom naturel $\widetilde G_E\mathbin{-}\widetilde G_F$\inserted{qui est un épimorphisme (\unsure{donc} surjectif)}, et de cette façon les $(G_E)$

\scanpage{186}
\mspage{II.34}
% archive sheet number: 185
% block: top, lines 1-2 | text
forment un système projectif de groupes algébriques, noté $\widetilde G=\pi^{\mathrm{alg}}=(\widetilde G_E)$. On définit de façon triviale

% block: upper third, lines 3-6 | diagram
\diagram{\begin{tikzcd}[ampersand replacement=\&,column sep=large]
\operatorname{Mod}(\pi^{\mathrm{alg}})\arrow[d,no head] \& \operatorname{Mod}_{\mathbb Q_\ell}(\pi)\arrow[l,maps to,"\simeq"']\\
\begin{gathered}\text{ensemble de}\\\text{représentations}\\\text{\uline{algébriques}}\\\text{de }\pi^{\mathrm{alg}}\end{gathered}\&
\end{tikzcd}}

% block: upper to middle third, lines 7-8 | text
qui est une équivalence de catégories. On aura d'ailleurs

% block: middle third, line 9 | display
\[
\pi^{\mathrm{alg}}/(\pi^{\mathrm{alg}})^0\simeq\pi
\]

% block: middle third, line 10 | text
ce qui signifie que, si $\pi^{\mathrm{alg}}=(\underset{\text{groupe alg}}{\widetilde G_\alpha})$, alors

% block: middle third, line 11 | display
\[
\pi\unsure{\simeq|\approx}(\widetilde G_\alpha/\widetilde G_\alpha^0)
\]

% block: middle third, lines 12-13 | text
(isom \unsure{canoniques|canonique}), où bien sûr on identifie les groupes finis ordinaires à des groupes alg. sur $\mathbb Q_\ell$.
\margin{\unsure{C algé.|C algri.|G alg.}}

% block: middle to lower third, lines 14-15 | text
\deleted{6.6} Au système projectif $\pi^{\mathrm{alg}}=(\widetilde G_\alpha)$ est associé une pro-Algèbre de Lie

% block: lower third, line 16 | display
\[
\deleted{\unsure{\mathfrak g}}\widetilde{\mathfrak g}=\operatorname{Lie}^{\mathrm{alg}}(\pi)^{\deleted{\unsure{\mathrm{an}}}}=\operatorname{Lie}(\widetilde G^{\deleted{\unsure{\mathrm{alg}}}})=(\operatorname{Lie}(\widetilde G_\alpha)),
\]

% block: lower third, line 17 | text
et on \deleted{\unsure{a}} a un foncteur évident

% block: lower third, lines 18-20 | diagram
\diagram{\begin{tikzcd}[ampersand replacement=\&,column sep=large]
\operatorname{Mod}_{\mathbb Q_\ell}(\pi)\arrow[r,phantom,"\simeq"] \& \operatorname{Mod}(\deleted{\pi^{\mathrm{alg}}}\widetilde G)\arrow[r] \& \operatorname{Mod}(\widetilde{\mathfrak g})\deleted{\unsure{\operatorname{Mod}(\mathfrak g)}}\arrow[d,no head]\\
\&\&\begin{gathered}\varinjlim\text{ des cat. de}\\\text{modules sur les }\widetilde{\mathfrak g}_\alpha\\\text{de dim fin. sur }\mathbb Q_\ell\end{gathered}
\end{tikzcd}}

% block: lower third, line 21 | text
\deleted{\unsure{on obtient encore}}\inserted{et}

% block: bottom, lines 22-26 | diagram
\diagram{\begin{tikzcd}[ampersand replacement=\&,column sep=large]
\varinjlim_{\unsure{U_i|H_i}}\operatorname{Mod}_{\mathbb Q_\ell}(U_i)\arrow[r,"\simeq"]\arrow[d,no head]
\&\operatorname{Mod}(\widetilde G^0)\arrow[r]\arrow[d,no head]\&\operatorname{Mod}(\widetilde{\mathfrak g})\\
\begin{gathered}\text{parcourt les}\\\text{s-groupes}\\\text{ouverts de }\unsure{\pi|U_i}\end{gathered}
\&\begin{gathered}\text{syst. projectif}\\\text{des composantes}\\\text{connexes }G_\alpha^0\\\text{des }G_\alpha\end{gathered}\&
\end{tikzcd}}

\scanpage{187}
\mspage{II.35}
% archive sheet number: 186
% block: upper third, lines 1-7 | text
La première flèche est une \uline{équivalence} \unsure{manifeste|naturelle} de catégories: le fait de restreindre \inserted{\unsure{des|les} représentations $\ell$-adiques} \unsure{mod. par ex.|mod. par gr.} à un s-groupe ouvert de $\pi$, revient à \deleted{ne} ne retenir de $\widetilde G=\pi^{\mathrm{alg}}$ que la \uline{composante connexe} $\widetilde G^0=(\widetilde G_\alpha^0)$. (N.B. Il \unsure{n'était|n'est} pas question de parler de composantes connexes dans le contexte analytique). Par contre, la flèche $\operatorname{Mod}(\widetilde G^0)\longrightarrow\operatorname{Mod}(\widetilde{\mathfrak g})$ est seulement \uline{pleinement fidèle}, mais pas \deleted{fidèle} \deleted{\unsure{à pri}} \unsure{ess.t} surjectif à priori.

% block: middle third, lines 8-13 | text
6.6. \deleted{\unsure{Nous|Alors}} Pour un groupe algébrique \inserted{affine} $\Gamma$ sur $\mathbb Q_\ell$, nous écrirons aussi $\Gamma$ au lieu de $\Gamma(\mathbb Q_\ell)$. Cela n'a pas d'inconvénient si les pts de $\Gamma/\Gamma^0$ sont rationnels sur $\mathbb Q_\ell$, donc $\Gamma(\mathbb Q_\ell)$ dense \unsure{dans|en} $\Gamma$, comme c'est \deleted{\unsure{donc}} le cas pour les groupes algébriques envisagés dans la section précédente.

% block: lower third, line 14 | text
Ceci dit, nous avons un hom canonique

% block: lower third, lines 15-16 | diagram
\diagram{\begin{tikzcd}[ampersand replacement=\&,column sep=large]
\text{(6.6.1)}\&\pi^{\mathrm{an}}\arrow[r]\arrow[d,equal]\&\pi^{\mathrm{alg}}\arrow[d,equal]\\
\&G\&\widetilde G
\end{tikzcd}}

% block: lower third, line 17 | text
qui s'explicite en écrivant

% block: lower third, line 18 | display
\[
\pi^{\mathrm{an}}\xrightarrow{\unsure{\sim}}\underset{E\in\operatorname{Ob}\operatorname{Mod}_{\mathbb Q_\ell}(\pi)}{\text{``}\varprojlim\text{''}}\unsure{G_E|G'_E}
\]

% block: bottom, lines 19-20 | text
or si $G_E$ est l'image de $\pi$ dans $\operatorname{Aut}_{\mathbb Q_\ell}(E)$, et

\scanpage{188}
\mspage{II.36}
% archive sheet number: 187
% block: top, line 1 | display
\[
\pi^{\mathrm{alg}}=\deleted{\unsure{\operatorname{Lie}_\ell|\varprojlim}}\text{``}\varprojlim\text{''}\widetilde G_E
\]

% block: upper third, lines 2-3 | text
où $\widetilde G_E$ est l'\uline{enveloppe algébrique} du groupe linéaire $G_E$.

% block: upper third, lines 4-6 | text
\uline{NB} En fait, la flèche $\unsure{\alpha}$ est un \uline{isom}, car tt groupe analytique compact sur $\mathbb Q_\ell$ a une représentation linéaire fidèle\ldots
\margin{\unsure{réf?}}

% block: middle third, lines 7-10 | text
Donc $\pi^{\mathrm{an}}\longrightarrow\pi^{\mathrm{alg}}$ est un isom de \deleted{\unsure{groupes}}\inserted{pro-groupes} analytiques ssi \inserted{tous} les $G_E$ sont \inserted{tous} des groupes \deleted{\unsure{linéaires}} algébriques. \inserted{Ce n'est sans doute jamais le cas, sauf si $\pi\simeq\pi^{\mathrm{an}}$ (donc $\pi\longleftarrow\pi^{\mathrm{alg}}$)}
\margin{?}

% block: middle third, lines 11-13 | deletion
\deleted{\unsure{Il suffit p.ex. que ce soient} leurs propres groupes dérivés --- en particulier qu'ils soient semi simples\ldots}

% block: middle third, line 14 | text
L'hom (6.6.1) induit un hom

% block: middle to lower third, line 15 | display
\[
\mathfrak g=\operatorname{Lie}^{\mathrm{an}}(\pi)\unsure{\longrightarrow|\text{---}}\widetilde{\mathfrak g}=\operatorname{Lie}^{\mathrm{alg}}(\pi).
\tag*{(6.6.2)}
\]

% block: lower third, lines 16-19 | text
C'est un isomorphisme si et seulement si les Algèbres de Lie $\mathfrak g_E=\operatorname{Lie}(G_E)$ sont algébriques, -- p.ex si ce sont leurs propres Algèbres dérivées, i.e. les $G_E$ sont \deleted{d'indice}\inserted{ouverts} dans les $\widetilde G_E$. Dans ce cas, \deleted{\unsure{l'hom}} foncteur

% block: lower third, lines 20-21 | display
\[
\varinjlim_{\substack{U_i\text{ ouvert}\\\text{de }\pi}}\operatorname{Mod}_{\mathbb Q_\ell}(U_i)\simeq\operatorname{Mod}(\widetilde{\mathfrak g})
\tag*{(6.6.3)}
\]

% block: lower third, line 22 | text
est évidemment une équivalence de catégories\ldots.

% block: bottom left margin | margin note
% slanted
\margin{laisser tomber remarques 6.7. et 6.8. dans \unsure{un texte|au net|au verso}}

% block: bottom, lines 23-24 | text
\uline{Remarque} 6.7. Si on interprète le premier \unsure{mb|mbr} en (6.6.3) comme $\operatorname{Mod}(\pi^{\mathrm{alg}})$, l'équivalence

\scanpage{189}
\mspage{II.37}
% archive sheet number: 188
% block: upper third, lines 1-9 | text
peut surprendre, car (\unsure{en admettant|mettant} que les $\widetilde G_E$ soient semi-simples disons) \deleted{\unsure{on ne voit pas qu'ils soient}} ils ne sont pas nécessairement simplement connexes (au sens de Chevalley) i.e. ils pourraient avoir (a priori\ldots) moins de représentations que leurs Algèbres de Lie. L'explication provient du fait que dans le système projectif des $\widetilde G_{E'}$ \unsure{majorant|dominant} $\widetilde G_E$, \unsure{on a|on trouve} aussi le ``revêtement simplement connexe'' de $\widetilde G_E$, au sens de Chevalley.

% block: middle third, lines 10-16 | text
\deleted{Pour}\inserted{\unsure{on}} s'en \unsure{assure|assurer} directement, \deleted{notons que,} \deleted{\unsure{soit} $\pi_E=\unsure{e|a}$}\inserted{soit} $G'$ le rev. simplement connexe de $\widetilde G_E^0$; et soit $\Gamma$ le noyau de $\deleted{G'\longrightarrow\widetilde G_E}G'(\mathbb Q_\ell)\unsure{\text{---}|\longrightarrow}\widetilde G_E^0(\mathbb Q_\ell)$. C'est un sous-groupe discret de $G'$, donc il existe un sous-groupe ouvert $U$ de $G'(\mathbb Q_\ell)$ tel que $U\longrightarrow\widetilde G_E^0(\mathbb Q_\ell)$ soit une \deleted{\unsure{isom}} application injective (\unsure{i.e. ouverte}). Si $E'$ est le module d'une représentation \inserted{alg} linéaire fidèle de $G'$, il induit une représentation \inserted{analytique} de $U$, donc une représentation d'un $\widetilde G_{E'}^0$\ldots

% block: middle left margin | margin note
% slanted
\margin{laisser tomber}

% block: lower third, lines 17-23 | text
6.8. Si $G$ est un groupe algébrique \deleted{affine}\inserted{\unsure{algébrique|quelconque|linéaire} (\unsure{opérant ds|plutôt de} $\unsure{E|G}$)} défini sur $\mathbb Q_\ell$, alors on peut considérer le groupe analytique $\deleted{\unsure{\pi}}=G(\mathbb Q_\ell)$; \deleted{et le groupe analytique profini \unsure{associé}}\inserted{\unsure{soit} $\pi$ un s-groupe ouvert, \uline{compact}}. Alors \deleted{\unsure{on peut considérer} $\pi$ \unsure{qui est}} le groupe pro-algébrique \inserted{$\pi^{\mathrm{alg}}$} associé est une \uline{extension} de $\pi=\deleted{\unsure{\widetilde\pi}}\unsure{\pi|\overline\pi}$ (considéré comme groupe pro-fini\inserted{-alg}) par

\scanpage{190}
\mspage{II.38}
% archive sheet number: 189
% block: upper third, lines 1-6 | text
un \deleted{\unsure{un}} groupe pro-algébrique connexe $(\widetilde G_{E'}^0)_{E'}$, \inserted{formé des}\deleted{\unsure{enveloppes}} \inserted{composantes connexes des}\unsure{algébriques|enveloppes} algébriques des représentations \deleted{\unsure{de de groupes}} $\ell$-adiques de $\pi$. Si p.ex. l'algèbre de Lie $\mathfrak g=\pi^{\mathrm{an}}=\operatorname{Lie}(G)$ est sa propre algèbre dérivée, alors dans toutes ses représentations $\ell$-adiques, l'image $G_{E'}$ de $\pi$ est ouverte dans $\widetilde G_{E'}$, donc $\widetilde G_{E'}^0\unsure{\text{---}|\longrightarrow}\widetilde G_E^0$ \deleted{induit un iso} est un isom local au sens $\ell$-adique, donc \deleted{\unsure{induit}} est \inserted{étale},

% block: upper right to middle, struck lines 7-9 | deletion
\deleted{\unsure{donc est isom si l'on suppose} $G^0$ \unsure{semi simple} et simplement connexe,}

% block: upper left to middle, beside cancelled text | diagram
\diagram{\begin{tikzcd}[ampersand replacement=\&,column sep=large]
\pi=G_{E'}\arrow[r,phantom,"{\unsure{\subset|\longleftarrow}}"]\arrow[d]\&\widetilde G_{E'}\arrow[d]\\
\pi=G_E\arrow[r,leftrightarrow]\&\widetilde G_E=G^1\quad\begin{gathered}\text{ss-groupe alg. ouvert}\\\text{de $G$ contenant $G^0$}\end{gathered}
\end{tikzcd}}

% block: middle third, lines 10-16 | text
Supposant pour simplifier $G^0$ \deleted{simple} $\frac12$ simple et simplement connexe, on trouve donc que l'enveloppe algébrique $G^1$ de $\pi$ (ég. à $G$ si $G$ est connexe) \unsure{apparait|apparaît} comme quotient d'une extension bien déterminée de $\pi$ (considéré comme groupe profini-algébrique) par $G'=\mathrm{rev.}$ \deleted{universel}\inserted{simplement connexe} de $G^0$; donc \unsure{on a|c.a.} un système projectif d'extensions de groupes algébriques

% block: middle left margin | margin note
% slanted
\margin{laisser tomber}

% block: lower third, line 17 | display
\[
1- G'\longrightarrow G_\alpha-\pi_\alpha-1
\]

% block: lower third, lines 18-20 | text
où $\pi_\alpha$ parcourt les quotients de $\pi$ par des ss-groupes \inserted{ouverts invariants}, $G^1$ étant lui-$=$ un des $G_\alpha$ \deleted{\unsure{$\Gamma$}}\inserted{\deleted{divisé} \unsure{par le} noyau de $G'\longrightarrow G^0$} ($\unsure{G^1|G'}/(G^1)^0$ étant évidemment un des quotients $\pi_\alpha$ de $\pi$\ldots).

% block: lower left margin, sideways note | margin note
% margin note written sideways
\margin{i.e. $\unsure{\xi|\zeta}$ est défini par l'extension $G^1$ de $G^1/G^0$ par $G^0$}

% block: lower third, line 21 | text
On en conclut un élément canonique

% block: lower third, line 22 | display
\[
\xi\in H^2(\pi,Z'(\mathbb Q_\ell))
\]

% block: bottom, lines 23-27 | text
où $Z'$ est le \deleted{noyau de}\inserted{centre de} $G'$\deleted{\unsure{\text{sur }G}}, l'annulation de $\xi$ étant nécessaire et suffisante pour que $\pi$ se remonte dans $G'$ en laissant fixe \unsure{par autom.} un Borel. L'image de $H^2(\pi,Z'(\mathbb Q_\ell))$ dans $H^2(\pi,Z(\mathbb Q_\ell))$ ($Z$ le centre de $G$) \deleted{\unsure{est finie et première}} est l'image de $\xi_0\in H^2(G^1/G^0,Z(\mathbb Q_\ell))$, \unsure{$\}$}

\scanpage{191}
\mspage{II.39}
% archive sheet number: 190
% block: top, lines 1-2 | text
6.9. \uline{Fonctorialité} Un hom de groupes profinis

% block: upper third, line 3 | display
\[
u:\pi\unsure{\text{---}|\longrightarrow}\pi'
\]

% block: upper third, line 4 | text
définit

% block: upper third, lines 5-6 | display
\[
\begin{aligned}
G=\pi^{\mathrm{an}}&\longrightarrow G'=\pi'^{\mathrm{an}},&\operatorname{Lie}(\deleted{\unsure{G}}\pi^{\mathrm{an}})=\mathfrak g&\text{---}\operatorname{Lie}(\pi'^{\mathrm{an}})=\mathfrak g'\\
\widetilde G=\pi^{\mathrm{alg}}&\longrightarrow\widetilde G'=\pi'^{\mathrm{alg}},&\operatorname{Lie}(\deleted{\unsure{G}}\pi^{\mathrm{alg}})=\widetilde{\mathfrak g}&\text{---}\operatorname{Lie}(\pi'^{\mathrm{alg}})=\widetilde{\mathfrak g}'
\end{aligned}
\]

% block: upper third, line 7 | text
avec commutativité dans

% block: upper to middle third, lines 8-9 | diagram
\diagram{\begin{tikzcd}[ampersand replacement=\&]
G=\pi^{\mathrm{an}}\arrow[r]\arrow[d,"u^{\mathrm{an}}"]\&\widetilde G=\pi^{\mathrm{alg}}\arrow[r]\arrow[d,"u^{\mathrm{alg}}"]\&\pi\arrow[d,"u"]\\
G'=\pi'^{\mathrm{an}}\arrow[r,no head]\&\widetilde G'=\pi'^{\mathrm{alg}}\arrow[r,no head]\&\pi'
\end{tikzcd}}

% block: middle third, lines 10-11 | diagram
\diagram{\begin{tikzcd}[ampersand replacement=\&,column sep=large]
\mathfrak g^{\deleted{\unsure{\mathrm{an}}}}=\operatorname{Lie}(\pi^{\mathrm{an}})\arrow[r,no head]\arrow[d]\&\mathfrak g'=\operatorname{Lie}(\pi'^{\mathrm{an}})\arrow[d]\\
\widetilde{\mathfrak g}=\operatorname{Lie}(\pi^{\mathrm{alg}})\arrow[r,no head]\&\widetilde{\mathfrak g}'=\operatorname{Lie}(\pi'^{\mathrm{alg}}).
\end{tikzcd}}

% block: middle third, line 12 | text
\unsure{Transitivités|Transitivité} \unsure{évidentes|évidente}.

% block: middle to lower third, lines 13-17 | text
Si $u:\pi\longrightarrow\pi'$ est un épimorphisme, il en est de même de $u^{\mathrm{an}}$ et $u^{\mathrm{alg}}$ (dans un sens évident). Si l'image de $u$ est d'indice fini, alors $u$ induit un \uline{épimorphisme} $G^0\text{---}G'^0$, et \unsure{des} épimorphismes $\mathfrak g\longrightarrow\mathfrak g'$, $\widetilde{\mathfrak g}\text{---}\widetilde{\mathfrak g}'$.

% block: lower third, lines 18-22 | text
6.10. Soit $\Pi=(\pi_{\unsure{\lambda|\alpha}})$ un système projectif \inserted{filtrant} de groupes profinis. En vertu de 6.9, il définit un système projectif \deleted{filtrant} de systèmes projectifs \deleted{filtrants} de groupes analytiques, groupes algébriques, Alg. de Lie et Alg. de Lie algébriques.

% block: bottom, lines 23-24 | display
\[
\begin{aligned}
G&=\Pi^{\mathrm{an}},&\operatorname{Lie}(\Pi^{\mathrm{an}})&=\mathfrak g\\
\widetilde G&=\Pi^{\mathrm{alg}},&\operatorname{Lie}(\Pi^{\mathrm{alg}})&=\widetilde{\mathfrak g}.
\end{aligned}
\]

\scanpage{192}
\mspage{II.40}
% archive sheet number: 191
% block: top, line 1 | text
Ce dernier donne lieu encore à des hom.

% block: upper third, lines 2-3 | display
\[
\begin{aligned}
\Pi^{\mathrm{an}}&\unsure{\text{---}|\longrightarrow}\Pi^{\mathrm{alg}}\text{---}\Pi\\
&\qquad\deleted{\unsure{\alpha}}
\end{aligned}
\]

% block: upper third, lines 4-5 | display
\[
\begin{aligned}
\text{et}\qquad\Pi&\simeq\Pi^{\mathrm{alg}}/(\Pi^{\mathrm{alg}})^0\\
\mathfrak g=\operatorname{Lie}(\Pi^{\mathrm{an}})&\text{---}\widetilde{\mathfrak g}=\operatorname{Lie}(\Pi^{\mathrm{alg}}).
\end{aligned}
\]

% block: upper third, lines 6-7 | text
Fonctorialité pour $\Pi$ variable comme pro-objet de la catégorie des groupes profinis.

% block: middle third, lines 8-10 | text
Attention, $\Pi^{\mathrm{an}}$ et \deleted{\unsure{de}} $\pi'^{\mathrm{an}}$ ne sont pas pareils, si $\pi'=\varprojlim\pi_\lambda$. \deleted{\unsure{pro-objet|pour}}\inserted{\unsure{D°|D'} pour $\Pi^{\mathrm{alg}}$, et les $\mathfrak g$, $\widetilde{\mathfrak g}$.} \deleted{\unsure{supposé|sauf} \unsure{par|pour} les morphismes de \unsure{transitivité|transition}}
% block: left margin beside Attention | drawing
\margin{\figure{\begin{tikzpicture}[x=1mm,y=1mm]
\draw (0,17) .. controls (11,20) and (21,17) .. (17,13)
.. controls (12,8) and (0,2) .. (3,0)
.. controls (6,-1) and (11,0) .. (13,1);
\end{tikzpicture}}}

% block: middle third, lines 11-16 | text
\uline{Exemple} 6.10.1. Les $\pi_\lambda$ finis, morphismes de transition surjectifs, $\pi'=\varprojlim\pi_\lambda$ groupe analytique de dim $>0$. Alors $\Pi^{\mathrm{an}}$ n'est autre que $(\pi_\lambda)$ lui-$=$, est formé de groupes discrets, mais $\pi^{\mathrm{an}}$ contient \inserted{\unsure{``}comme facteur''} $\pi'$! [\inserted{Ici} $\Pi^{\mathrm{an}}$ \deleted{plus} est ``plus petit'' que $\pi'^{\mathrm{an}}$!].

% block: lower third, lines 17-23 | text
\uline{Exemple} 6.10.2. Les $\pi_\lambda$ les sous-groupes d'indice fini d'un $\pi$. Alors $\pi'=\{e\}$, donc $\pi'^{\mathrm{an}}$ et $\pi'^{\mathrm{alg}}$ sont les pro-groupes unités. Mais $\Pi^{\mathrm{an}}$ est le système des ss-groupes \deleted{d'indice fini}\inserted{ouverts} des $G_\alpha$ définissant $\pi^{\mathrm{an}}$, et $\Pi^{\mathrm{alg}}$ n'est autre que $\widetilde G^0$, où $\widetilde G=\pi^{\mathrm{alg}}$! \deleted{Donc} [Ici $\Pi^{\mathrm{an}}$ est ``plus grand'' que $\pi'^{\mathrm{an}}$, qui est trivial!]

\scanpage{193}
\mspage{II.41}
% archive sheet number: 192
% block: top, 6.11 | display
\[
6.11.\qquad \operatorname{Mod}_{\mathbb{Q}_\ell}(\Pi^{\mathrm{an}})\simeq
\operatorname{Mod}(\Pi^{\mathrm{alg}})\simeq
\varinjlim_\lambda\operatorname{Mod}_{\mathbb{Q}_\ell}(\pi_\lambda)
\]
% block: upper third, 6.12 | text
6.12. Un cas \unsure{important} est celui où dans le \deleted{système}
pro-objet $\Pi=(\pi_\lambda)$, les morphismes de transition sont
\unsure{quasi-surjectifs} (i.e. image d'indice fini). Alors il en est de même
dans $\Pi^{\mathrm{an}}$ et $\Pi^{\mathrm{alg}}$ ; (en particulier
$(\Pi^{\mathrm{alg}})^0$ est un pro-objet \uline{strict}). De plus, les syst.
projectifs $\mathfrak{g}=\operatorname{Lie}(\Pi^{\mathrm{an}})$ et
$\widetilde{\mathfrak{g}}=\operatorname{Lie}(\Pi^{\mathrm{alg}})$ sont aussi
stricts \unsure{....|...}

% block: middle third, 6.13 | text
6.13. Soit $X=\varprojlim_\lambda X_\lambda$ un schéma, limite projective
d'un système projectif filtrant \inserted{$\unsure{\mathscr{X}|\mathcal{X}}$}
de schémas, à morphismes de transition affines. Soit $\unsure{\xi|\zeta}$ un pt géom.
de $X$, $\xi_\lambda$ son image dans $X_\lambda$. Considérons
% block: middle, fundamental pro-group | display
\[
\Pi=\pi_1(\mathscr{X},\xi)=(\pi_1(X_\lambda,\xi))_\lambda,
\]
% block: lower third, representations | text
système projectif de groupes profinis. Les représentations
$\ell$-adiques de ce \deleted{système} pro-groupe top. forment une
catégorie équivalente à la catégorie
$\varinjlim_\lambda(\inserted{\mathbb{Q}_\ell\text{-}}\text{Faisceaux
\deleted{constants} tordus sur }X_\lambda)$, d'où leur intérêt. Cette
catégorie s'explicite en termes de $\Pi^{\mathrm{an}}$ ou $\Pi^{\mathrm{alg}}$
comme on a vu plus haut, et peut aussi (partiellement) s'étudier à l'aide
de $\operatorname{Lie}(\Pi^{\mathrm{an}})$ et
$\operatorname{Lie}(\Pi^{\mathrm{alg}})$

\scanpage{194}
\mspage{II.42}
% archive sheet number: 193
% block: upper third, repeated 6.13 | text
6.13. Le cas le plus intéressant est celui où les $X_\lambda$ sont de type
\deleted{\unsure{de $X_\lambda$}} fini sur $\mathbb{Z}$, ou sur un corps
premier, \deleted{Alors} On sait que le foncteur
\inserted{$\unsure{\mathscr{X}|\mathcal{X}}\mathbin{-}X$}
\unsure{limite} des pro-objets dans la catégorie des schémas de type fini
sur $\unsure{Z|\mathbb Z}$, dans celle des schémas, est pleinement fidèle, i.e.
$\mathscr{X}$ est déterminé par la connaissance de $X$.

% block: middle, example | text
\uline{Ex} $X$ de présentation finie sur un schéma affine, p.ex. sur un
corps. En particulier $X$ affine, p.ex. \unsure{$\alpha$|$X$|$=$} spectre d'un
corps, p.ex. d'un corps alg. clos (sic)

% block: middle, arithmetic group name | text
Dans ce cas, $\Pi$ mérite d'être \unsure{appellé|appelé} le
``\inserted{pro}\uline{groupe fondamental arithmétique de} $X$'',
\inserted{notation $\unsure{\pi_1^{\mathrm{arith}}|\Pi_1^{\mathrm{arith}}}(X,\unsure{\xi|\zeta})$}
(en se rappelant que c'est un \deleted{groupe} pro-groupe topologique.
On a un isom. canonique
% block: lower middle, canonical isomorphism | display
\[
\pi_1(X,\xi)\xrightarrow{\sim}\varprojlim\Pi
=\varprojlim_\lambda\pi_1(X_\lambda,\xi),
\]
% block: lower third, limitation | text
mais on \unsure{se} \unsure{rappellera} que la connaissance de $\pi_1(X,\xi)$ ne
détermine pas celle \deleted{de} \inserted{pro-groupe}
\deleted{\unsure{du}} $\Pi=\Pi_1(X,\xi)$.
Par exemple, si $X$ est le spectre d'un corps alg. clos,
$\pi_1(X,\xi)=e$, tandis que $\Pi$ n'est alors jamais trivial, comme
nous allons voir.

% block: foot, continuation | text
En vertu de 6.10, $\pi^{\mathrm{arith}}(X,\xi)$
donne

\scanpage{195}
\mspage{II.43}
% archive sheet number: 194
% block: top, continuation of 6.13 | text
naissance à des groupes proanalytiques et proalgébriques
% block: top, completions | display
\[
\unsure{\pi_1^{\mathrm{arith}}|\Pi_1^{\mathrm{arith}}}(X,\unsure{\xi|\zeta})^{\mathrm{an}},\quad
\pi_1^{\mathrm{arith}}(X,\xi)^{\mathrm{alg}}
\]
% block: upper third, Lie algebras | text
et à des proalgèbres de Lie correspondantes.

% block: upper third, 6.14 | text
6.14. Supposons $X$ intègre et normal, alors on peut supposer les
$X_\lambda$ intègres, normaux (quitte à \deleted{les} remplacer $X_\lambda$
par le normalisé de l'image schématique fermée de $X$ dans $X_\lambda$
$\ldots$). \deleted{\unsure{Les}}
Comme \unsure{pour} tt morphisme dominant de t.f. de schémas normaux,
$X_\mu\mathbin{-}X_\lambda$ induit un \deleted{quasi} hom.
\uline{quasi-surjectif}
$\pi_1(X_\mu,\xi_\mu)\mathbin{-}\pi_1(X_\lambda,\xi_\lambda)$,
i.e. on est dans le cas de 6.12. En particulier, les \deleted{syst.} hom.
de transition dans \deleted{$\unsure{\Pi}$}
$G=\pi_1^{\mathrm{arith}}(X,\xi)^{\mathrm{an}}$ et
$\widetilde{G}=\pi_1^{\mathrm{arith}}(X,\xi)^{\mathrm{alg}}$ sont des
quasi-épimorphismes, et les hom. de transition dans $\mathfrak{g}$ et
$\widetilde{\mathfrak{g}}$ sont des épimorphismes (i.e.
$\mathfrak{g},\unsure{\widehat{\mathfrak{g}}|\widetilde{\mathfrak{g}}}$
sont des pro-objets \uline{stricts}).

% block: lower third, 6.15 and interlinear corrections | text
6.15. Notons que \deleted{si $\ell$ est inversible sur $X$.}
\inserted{si $\ell$ est inversible sur $X$ (\unsure{seul cas intéressant}
$\ldots$) alors}
$\Pi\ (=\pi_1^{\mathrm{arith}}(X,\xi))$
n'est \deleted{\unsure{jamais}} ``discret'', i.e.
\deleted{\unsure{il y a d}} $\dim\mathfrak{g}>0$,
\deleted{$\unsure{\dim\mathfrak{g}}$} ou encore
$\dim\widetilde{\mathfrak{g}}>0$. (En fait, on verra que ces dimensions
sont \uline{tjs} infinies).

% block: foot, argument beginning | text
En effet, il suffit de prouver que si \deleted{et $\ell$ y est inversible}
$X$ est de type fini sur $\operatorname{Spec}\mathbb{Z}$, et muni d'un
pt géom. $\xi$, \inserted{et si $\ell$ y est inversible,} alors,
\deleted{$\unsure{\Pi_1(X,|\pi_1(X,}$} \unsure{posons}
$\pi=\pi_1(X,\xi)$, $\pi^{\mathrm{alg}}$ n'est jamais discret. Or
